% Pursuing Stacks (cotes 134-2 to 134-8): archive cote 134-3, site page 74 (PDF page 75), typescript page 74.
% typescript. Draft compared: ps2.tex:3821-3890.
% AI-transcribed by Claude Opus 5.5 (claude-opus-5-5, adaptive thinking, effort high), Message Batches API, batch ps1 (msgbatch_01Rv5hCaNTtUt4JHx1yqoHPL), 2026-11-12; not yet checked by eye.
\scanpage{134-3/74}
\mspage{74}
and h'\inserted{$:M\longrightarrow M'$} induce functors between $W'^{-1}M'$ and $W^{-1}M$, which are quasi-inverse
to each other - hence these two categories are equivalent. In case
(M,W) is a modelizer, this implies that (M',W') is a modelizer too and the
\inserted{inclusion} functor\xout{s} g, as well as h', are model\xout{s} preserving.

We can apply this remark to the modelizer (Cat), and to the functor
\[ i_A\, i^{*} : (Cat) \longrightarrow (Cat) \]
defined by any functor i: $A\longrightarrow$ (Cat) satisfying the conditions of the
``key theorem'' (\S. \over{5}{3}7), for instance the functor $i_A\hand{|}A$ : $a\longmapsto A_{/a}$
(where A is any test category). We get the following
\uline{Proposition} Let M be any full subcategory of (Cat), assume there exists
a test-category A such that for any F in $A\hat{}$, the category $A_{/F}$ belongs to
M (i.e. the functor $i_A$ : $A\hat{}\longrightarrow$ (Cat) , $F\longmapsto A_{/F}$ , factors through M). \xout{Then}
Let $W_M$ be the set of weak equivalences in M . Then $(M,W_M)$ is a \xout{test cate-}%
\xout{gory} \inserted{modelizer} and the inclusion functor $M\to$ (Cat) is model preserving.

The same of course will be true for any full subcategory of (Cat) con-
taining M - which makes an impressive bunch of modelizers indeed ! When
\inserted{the test category} A is given, \over{O}{o}ne natural choice for M is to take all categories \inserted{C} which are
``locally isomorphic to A '' , namely such that for any x in C, the indu-
ced category $C_{/x}$ be isomorphic to a category of the type $A_{/a}$ , with a in
A .

It would be tempting to apply this result to the full subcategory
(Ord) of (Cat) - but for this to be feasible, would mean exactly that
there exists a test-category A defined by an ordered set\inserted{\deleted{(or at least ``locally ordered'')}}. To see whether
there exists indeed such an ordered set looks like a rather interesting
question - maybe it would give rise to \inserted{algebraic} models for homotopy types, sim-
pler than those used so far, namely simplicial and cubical complexes and
multicomplexes. It is interesting to note that if such a test category
should exist, it will \uline{not} be in any case a strict test category. Indeed,
the topos $A\hat{}$ associated to an ordered set A can be viewed also, as we saw
before (\S.\unsure{22}), as associated to \inserted{a suitable} \deleted{the} topological space (namely A endowed
with a suitable topology, the open sets being just the``cribles'' in A).
But we have seen that the topos associated to a topological space \uline{cannot}
be \uline{strictly} modelizing (\inserted{prop.2,} cor.2 \over{o}{i}n \deleted{page 53}\inserted{section 35}). \hand{$\rfloor$} This remark confirms the fee-
ling that it was worth while emphasizing the notion of a test category
\margin{X} (just T 1 + T 3 ) by a simple and striking name as I finally did, rather
than bury it behind the notion I now call a strict test category, which
is considerably more stringent and, moreover, more``rigid''. For instance,
it is not stable under localization $A_{/a}$ , whereas the notion of a test cate-
gory is - indeed, for any aspheric I in $A\hat{}$, $A_{/I}$ is still a test categor\hand{y}
